% Pregunta 782 (2005-septiembre-OEP2005, n.º 1): mercado de dinero keynesiano. Al aumentar la oferta
% monetaria (M0 -> M1) el tipo baja; cuando la oferta corta el tramo horizontal de la demanda
% (trampa de la liquidez, M2 y M3) el tipo ya no baja más (b correcta, a falsa).
% Autor: Víctor Gutiérrez Marcos
\documentclass[tikz,border=6pt]{standalone}
\input{tikz-tcee}
\begin{document}
\begin{tikzpicture}[tcee]
  \ejes{$M/P$}{$i$}{6.6}{5}
  % Demanda de dinero: decreciente y convexa, tangente al suelo i_min en M/P = 3.6
  \draw[curva1,domain=0.6:3.6,samples=60] plot (\x,{0.8+3.2*((3.6-\x)/3)^2}) -- (6.2,0.8)
    node[etiqueta,above,text=tceemorado] {$L$};
  % Ofertas monetarias
  \draw[curva2] (1.4,0) -- (1.4,4.3) node[etiqueta,above] {$M_0$};
  \draw[curvanueva,desplazada] (2.4,0) -- (2.4,4.3) node[etiqueta,above] {$M_1$};
  \draw[curvanueva,desplazada] (4.3,0) -- (4.3,4.3) node[etiqueta,above] {$M_2$};
  \draw[curvanueva,desplazada] (5.3,0) -- (5.3,4.3) node[etiqueta,above] {$M_3$};
  \draw[flecha] (1.55,3.9) -- (2.25,3.9);
  \draw[flecha] (2.55,3.9) -- (4.15,3.9);
  \draw[flecha] (4.45,3.9) -- (5.15,3.9);
  % Equilibrios y tipos
  \draw[guia] (0,2.521) -- (1.4,2.521);
  \draw[guia] (0,1.312) -- (2.4,1.312);
  \node[etiqueta,left] at (0,2.521) {$i_0$};
  \node[etiqueta,left] at (0,1.312) {$i_1$};
  \node[etiqueta,left] at (0,0.8) {$i_{\min}$};
  \draw[guia] (0,0.8) -- (3.6,0.8);
  \node[punto] at (1.4,2.521) {};
  \node[punto] at (2.4,1.312) {};
  \node[punto] at (4.3,0.8) {};
  \node[punto] at (5.3,0.8) {};
  % Rótulo del tramo horizontal (con flecha hacia el suelo)
  \draw[->,thin,tceemorado] (3.9,1.45) -- (4.6,0.95);
  \node[etiqueta,font=\footnotesize,text=tceemorado,align=center] at (3.35,2.05) {trampa de la\\liquidez: $i$ ya\\no baja más};
\end{tikzpicture}
\end{document}
